\documentclass[11pt]{article}
\usepackage[T1]{fontenc}
\usepackage{lmodern}
\usepackage{amsmath,amssymb,amsthm,mathtools}
\usepackage[a4paper,margin=28mm]{geometry}
\usepackage[hidelinks]{hyperref}
\newtheorem{theorem}{Theorem}[section]
\newtheorem{proposition}[theorem]{Proposition}
\newtheorem{lemma}[theorem]{Lemma}
\newtheorem{corollary}[theorem]{Corollary}
\theoremstyle{remark}
\newtheorem{remark}[theorem]{Remark}
\title{Lower bounds for polynomial definitions of rings of integers}
\author{Research supplement}
\date{6 November 2026}
\begin{document}
\maketitle
\begin{abstract}
We record unconditional restrictions on polynomial existential definitions
of rings of integers in number fields. Besides the bounds of total degree
at least four and at least two existential witnesses, we prove that a
system of quadratic equations in at most four total affine coordinates
cannot have an infinite input projection contained in the ring of integers.
Consequently a quartic defined as a positive sum of squares of quadratics
requires at least five total coordinates. A globally nonnegative polynomial
in three rational coordinates whose rational projection is an infinite
subset of the integers has degree at least eight. The unrestricted
sign-changing quartic surface case is distinguished from these restricted
results. We also prove that a finite union of equations jointly quadratic
in two witnesses, with unrestricted input degree, either admits a generic
rational section or misses infinitely many members of every integral
arithmetic progression. Such equations cannot define the ring of integers
or an integral superset with infinite complement. The arguments do not
assume a proposed Diophantine definition.
\end{abstract}
\tableofcontents
% Includable section. Requires amsmath, amssymb, amsthm, hyperref.
% Bibliography entries are supplied in lower_bounds.bib.
\section{Unconditional lower bounds}
\label{sec:lower-bounds}

The lower bounds in this section do not assume the correctness of any proposed
Diophantine definition.  In fact, they hold over every number field.
Let $K$ be a number field and let $\mathcal O_K$ be its ring of integers.
For a polynomial
\[
 F(T,Y_1,\ldots,Y_r)\in K[T,Y_1,\ldots,Y_r]
\]
write
\[
 S_F=\{t\in K:\text{there exist }y_1,\ldots,y_r\in K
                    \text{ with }F(t,y_1,\ldots,y_r)=0\}.
\]
The variable $T$ is the input; the $r$ variables $Y_i$ are the witnesses.
The degree is the total degree, including $T$.

\begin{theorem}\label{thm:lower-bounds}
If $S_F=\mathcal O_K$, then
\[
                    r\geq 2\qquad\text{and}\qquad\deg F\geq 4.
\]
In particular, a Diophantine definition of $\mathbb Z$ in $\mathbb Q$
by one polynomial needs at least three total variables and has degree at least
four.  Thus a valid quartic construction has the smallest possible degree,
regardless of its number of variables.
\end{theorem}

We prove the two assertions separately.  The variable bound uses the integral
form of Hilbert irreducibility.  The cubic obstruction uses Koll\'ar's theorem
on unirational cubic hypersurfaces; the treatment of cones and reducible
polynomials is included explicitly.

\subsection{One witness and Hilbert irreducibility}

We use the following standard form of Hilbert irreducibility.  Given finitely
many irreducible polynomials $f_i(T,Y)\in K(T)[Y]$, each of positive degree in
$Y$, and finitely many forbidden values of $T$, there are infinitely many
$a\in\mathcal O_K$ outside those values for which all the specializations are
defined, preserve their $Y$-degrees, and remain irreducible in $K[Y]$.
See \cite[Chapter~4, Section~4]{lbSchinzel2000}; the integral number-field
form is also recalled in \cite[Introduction]{lbMuller2001}.
The restriction to integral specializations is essential here: the
Hilbertianity of $K$ alone, without that strengthening, is not the assertion
being applied.

\begin{proposition}\label{prop:one-witness}
Let $G\in K[T,Y]$.  If $S_G$ contains all but finitely many elements of
$\mathcal O_K$, then $S_G$ contains all but finitely many elements of $K$.
\end{proposition}
\begin{proof}
The conclusion is immediate if $G=0$.  If $G\neq0$ and has degree zero in
$Y$, then $S_G$ is finite, contrary to the hypothesis.  Otherwise factor in
$K(T)[Y]$:
\[
                G=c(T)\prod_{i=1}^{s}f_i(T,Y)^{e_i},
\]
where $c(T)\neq0$, the $e_i$ are positive integers, and the $f_i$ are monic
irreducible polynomials of positive $Y$-degree.
If some $f_i$ is linear, it has a root $u(T)\in K(T)$.  The identity
$G(T,u(T))=0$ then specializes for all but finitely many $t\in K$,
proving the conclusion.

If every $f_i$ has $Y$-degree at least two, integral Hilbert irreducibility
provides infinitely many $a\in\mathcal O_K$ at which $c(a)\neq0$ and all
$f_i(a,Y)$ remain irreducible of degree at least two.  None has a root in
$K$, so $G(a,Y)$ has no root in $K$.  This contradicts the hypothesis.
\end{proof}

Since $K\setminus\mathcal O_K$ is infinite, the proposition rules out one
witness.  For example, $1/n\notin\mathcal O_K$ for every integer $n\ge2$,
because a rational algebraic integer is an integer.  With no witness a
nonzero polynomial in $T$ has only finitely many roots, whereas the zero
polynomial defines all of $K$.  This proves the variable assertion of
Theorem~\ref{thm:lower-bounds}.

\subsection{A topological obstruction to rational parametrizations}

Let $r_1$ and $r_2$ denote the numbers of real embeddings and conjugate pairs
of complex embeddings of $K$.  Put
\[
 K_{\infty}=\mathbb R^{r_1}\times\mathbb C^{r_2},
\]
viewed as a real vector space.  The diagonal embedding
$\iota:K\longrightarrow K_{\infty}$ has dense image, while
$\iota(\mathcal O_K)$ is a discrete lattice.  Indeed, an integral basis of
$K$ identifies these images with the images of $\mathbb Q^{[K:\mathbb Q]}$
and $\mathbb Z^{[K:\mathbb Q]}$ under an invertible real linear map.

\begin{lemma}\label{lem:rational-function-integral}
Let $R\in K(U_1,\ldots,U_m)$ be nonconstant, and let $U$ be a nonempty
Zariski-open subset of $\mathbb A^m_K$ on which $R$ is defined.
Then $R(U(K))$ is not contained in $\mathcal O_K$.
\end{lemma}
\begin{proof}
Assume to the contrary that all these values lie in $\mathcal O_K$.
After shrinking $U$ we may take $U$ to be a principal open subset, defined by
one nonzero polynomial.  Apply every archimedean embedding of $K$ to its
coefficients and to those of $R$.  This gives a continuous map, on a nonempty
Euclidean open subset of $K_{\infty}^{m}$,
\[
 R_{\infty}:\prod_{v\mid\infty}K_v^m
       \longrightarrow\prod_{v\mid\infty}K_v=K_{\infty}.
\]
Choose a small connected open ball $B$ in its domain.  The diagonal image of
$K^m$ is dense in $B$.  By hypothesis its image under $R_{\infty}$ lies in
the closed discrete set $\iota(\mathcal O_K)$; by continuity the same is true
on all of $B$.  A continuous map from a connected set to a discrete set is
constant.  Each archimedean component of $R$ is therefore constant on an
open set, so its numerator minus that constant times its denominator is the
zero polynomial.  This forces $R$ to be constant, a contradiction.
\end{proof}

For $K=\mathbb Q$, this argument is simply the density of $\mathbb Q$ in
$\mathbb R$ and the discreteness of $\mathbb Z$.  For a general number field
one must use all archimedean embeddings together; $\mathcal O_K$ need not be
discrete in one selected embedding.

An integral variety $X/K$ is called \emph{unirational over $K$} if some
projective space admits a dominant rational map to $X$, defined over $K$.

\begin{corollary}\label{cor:unirational-integral}
If $X$ is unirational over $K$ and $f\in K(X)$ is nonconstant, then the values
of $f$ at the $K$-points where it is defined cannot all belong to
$\mathcal O_K$.
\end{corollary}
\begin{proof}
Compose $f$ with a dominant rational parametrization of $X$, and restrict
the source to a suitable nonempty open subset of affine space.  Dominance
makes the resulting rational function nonconstant.  Apply
Lemma~\ref{lem:rational-function-integral}.
\end{proof}

\subsection{Quadratic polynomials}

\begin{lemma}\label{lem:quadratic-integral}
If $\deg F\le2$ and $S_F\subseteq\mathcal O_K$, then $S_F$ has at most two
elements.
\end{lemma}
\begin{proof}
Suppose that $a,b,c\in K^{r+1}$ are zeros of $F$ with three distinct input
coordinates.  If these points are collinear, the restriction of $F$ to their
line has degree at most two and three distinct roots.  Thus the entire
line is contained in $F=0$.  The input is nonconstant on that line and takes
all values in $K$, a contradiction.

Otherwise the three points span an affine plane $L$ over $K$.  The input
coordinate is a nonconstant affine linear function on $L$.  If
$F|_L=0$, a nonvertical line in $L$ again gives a contradiction.  If
$F|_L$ factors into affine linear factors over $K$, one of its at most two
lines contains two of the selected points, hence has nonconstant input;
this gives the same contradiction.  A nonzero linear polynomial cannot
vanish at three noncollinear points.

The remaining case is an irreducible quadratic polynomial on $L$.
It defines an irreducible conic with a smooth $K$-point: a singular conic
has at most one singular point, so the three distinct points cannot all be
singular.  Projection from a smooth $K$-point parametrizes the conic by
rational functions over $K$.  The input is nonconstant on this conic,
for otherwise an irreducible quadratic curve would be contained in an
affine line.  Its composition with the parametrization is a nonconstant
rational function taking only values in $\mathcal O_K$, contrary to
Lemma~\ref{lem:rational-function-integral}.
\end{proof}

\subsection{Cubic cones and the sharp degree obstruction}

We use the following published theorem, in its characteristic-zero case.
\begin{theorem}[Koll\'ar]\label{thm:kollar-input}
Let $K$ be a perfect field and let $X\subseteq\mathbb P^{n+1}_K$ be a
geometrically integral cubic hypersurface of dimension $n\ge2$ which is not a cone.
If $X(K)\neq\varnothing$, then $X$ is unirational over $K$.
\end{theorem}
This is the geometrically integral case of Theorem~2 of
\cite{lbKollar2002}; the theorem there is stated for irreducible cubics.
We handle geometric reducibility separately below, so no convention about
\emph{irreducible} is needed.  Its cone exclusion must not be omitted: a cone over a smooth plane cubic need not be unirational.

For a homogeneous cubic $H$ on a finite-dimensional $K$-vector space $E$,
define its vertex subspace by
\[
 W=\{v\in E:D_vH=0\text{ as a polynomial on }E\},
 \qquad D_vH(x)=\left.\frac{d}{ds}H(x+sv)\right|_{s=0}.
\]
The equations defining $W$ are linear over $K$.  In characteristic zero,
$D_vH=0$ is equivalent to $H(x+sv)=H(x)$ for all $x,s$.
Choosing a basis of $W$ and extending it to a basis of $E$ therefore makes
$H$ independent of the coordinates along $W$.
Consequently $H$ descends to a homogeneous cubic $\overline H$ on $E/W$.
By maximality of $W$, this quotient cubic has no nonzero vertex direction
and is not a cone.  These statements remain true after extending scalars:
the kernel of a matrix over $K$ commutes with field extension.

\begin{proof}[Proof of the degree assertion in Theorem~\ref{thm:lower-bounds}]
Suppose $S_F=\mathcal O_K$ and $\deg F\le3$.
The polynomial $F$ is neither zero nor a nonzero constant.
If $F$ is reducible, each of its irreducible factors has degree at most two.
In that case, the projection is the union of the
projections of finitely many linear or quadratic factors.  Every such
projection is contained in $\mathcal O_K$, hence finite by
Lemma~\ref{lem:quadratic-integral}.  This cannot give $\mathcal O_K$.
The same lemma excludes $\deg F\le2$.
We may therefore suppose that $F$ is irreducible of degree three.

Homogenize $F$ to an irreducible cubic
\[
 H(X_0,X_1,\ldots,X_{r+1})
      =X_0^3 F(X_1/X_0,\ldots,X_{r+1}/X_0).
\]
Here $X_1/X_0$ is the input function.
We first exclude geometric reducibility.  Since the characteristic is zero,
$H$ has no repeated geometric components.  Irreducibility over $K$ makes
the Galois group act transitively on them.  Their degrees are equal and
sum to three; thus if $H$ is geometrically reducible, it is a product of
three conjugate linear forms.  A $K$-rational zero belongs to every one of
these hyperplanes, because it is fixed by Galois.  Consequently the entire
line joining any two $K$-rational zeros lies in all three hyperplanes.
Taking affine zeros with inputs $0$ and $1$ would produce every input in
$K$, a contradiction.  We may therefore assume $H$ is geometrically
integral.

Let $W$ be its vertex subspace.
For any $v\in W$, any affine rational zero $x$ normalized by $x_0=1$,
and any $s\in K$ with $1+s v_0\neq0$, the point $x+sv$ is a zero of $H$.
Its input is
\[
                       \frac{x_1+s v_1}{1+s v_0}.
\]
If $v_1-x_1v_0\neq0$, this is a nonconstant M\"obius function of $s$;
its image omits at most one element of $K$.  It cannot be contained in
$\mathcal O_K$, because $K\setminus\mathcal O_K$ is infinite.
Hence $v_1=x_1v_0$ at every such point.  Apply this to zeros with input
$0$ and input $1$, whose existence follows from $S_F=\mathcal O_K$.
It follows that
\[
                         v_0=v_1=0\qquad(v\in W).
\]

Thus both $X_0$ and $X_1$ descend to independent linear forms on $E/W$.
Choose coordinates on the quotient beginning with these forms.  The
irreducible quotient cubic $\overline H$ is not a cone, and its affine
chart $X_0=1$ has exactly the same input projection as $H$.
Indeed, every rational point on the quotient lifts through a $K$-linear
splitting of $E\to E/W$, preserving $X_0$ and $X_1$, and $H$ is constant
along the fibers.

If the projective quotient hypersurface has dimension at least two,
Theorem~\ref{thm:kollar-input} makes it unirational, since it has a
$K$-point with $X_0=1$.  Its input function $X_1/X_0$ is nonconstant,
as its image contains all of the infinite set $\mathcal O_K$.
Corollary~\ref{cor:unirational-integral} now gives a nonintegral input,
a contradiction.

If the quotient has dimension one, its affine chart is defined by a
polynomial in the input and just one further coordinate.  Its projection
cannot be $\mathcal O_K$ by Proposition~\ref{prop:one-witness}.
If the quotient has dimension zero, the affine chart has only finitely
many points and likewise cannot project onto $\mathcal O_K$.
These cases exhaust the possibilities, proving $\deg F\ge4$.
\end{proof}

\begin{remark}[Scope of the optimality statement]
The conclusion concerns one polynomial equation with all witnesses in the
field, and total degree in all variables.  It does not prohibit a system
of quadratic equations: a sum of their squares over $\mathbb Q$ is quartic.
It also does not determine the smallest possible number of witnesses at
degree four, or the smallest degree among definitions attaining the
unknown minimum number of witnesses.  The unconditional numerical
bounds established here are two witnesses and degree four; no example
meeting both bounds is asserted.
\end{remark}

\subsection{A stronger one-witness obstruction}

For completeness there is a useful strengthening which uses Siegel's
theorem rather than just Hilbert irreducibility.
In its rational-function form, Siegel's theorem says that if $C/K$ is a
smooth projective geometrically integral curve of positive genus and
$f\in K(C)$ is nonconstant, then
\[
              \{P\in C(K):f(P)\text{ is defined and belongs to }
                       \mathcal O_K\}
\]
is finite.  One can derive this from the usual theorem on affine integral
points by taking $C$ minus the poles of $f$: the map to $\mathbb A^1$
defined by $f$ is finite, and after inverting finitely many primes its
coordinate functions are integral over $\mathcal O_{K,S}[f]$.
This reduction is also used in
\cite[Section~2.1]{lbMuller2001}; see \cite[Chapter~8]{lbLang1983}.

\begin{proposition}\label{prop:one-witness-finite}
If $G\in K[T,Y]$ satisfies $S_G\subseteq\mathcal O_K$, then $S_G$ is finite.
Thus no infinite subset of $\mathcal O_K$ is the projection of the
$K$-zeros of a single polynomial with just one witness.
\end{proposition}
\begin{proof}
Factor $G$ and treat its finitely many irreducible factors separately.
Factors depending only on $T$ give finite projections.  For another factor,
let $C$ be its plane curve.  If $C$ is not geometrically irreducible, its
$K$-points lie in the intersections of distinct geometric components,
a finite set.  Otherwise pass to its smooth projective normalization
$\widetilde C$.  Away from the finitely many singular points of $C$,
rational points lift to $\widetilde C(K)$, and the input defines a
nonconstant rational function $t$ on $\widetilde C$.

If $\widetilde C$ has positive genus, Siegel's theorem gives only finitely
many points for which $t$ is integral.  If its genus is zero and it has
any rational point, then $\widetilde C\simeq\mathbb P^1_K$.
The input then becomes a nonconstant rational function on
$\mathbb P^1_K$, and on a nonempty open subset all its values would be
integral, contradicting Lemma~\ref{lem:rational-function-integral}.
If it has no rational point, only the finitely many previously excluded
singular points remain.  Each case gives a finite projection.
\end{proof}

\begin{corollary}\label{cor:cubic-finite}
Let $F\in K[T,Y_1,\ldots,Y_r]$ have total degree at most three.
If $S_F\subseteq\mathcal O_K$, then $S_F$ is finite.
\end{corollary}
\begin{proof}
For an infinite projection the proof of
Theorem~\ref{thm:lower-bounds} applies unchanged: choose any two distinct
input values in the geometric-reducibility argument and in the vertex
elimination.  In the curve case use
Proposition~\ref{prop:one-witness-finite} instead of
Proposition~\ref{prop:one-witness}.  The other cases already give a
contradiction to infinitude.
\end{proof}

% Includable research supplement. Theorem environments and bibliography are
% provided by paper.tex. All variable counts include the input unless stated.

\section{Stronger obstructions for small polynomial presentations}
\label{p2lb:scope}

The unrestricted bounds established in the first report are two existential
witnesses and total degree four.  The arguments below give stronger bounds
for important classes of presentations.  They do not prove that every
definition needs more than two witnesses.  In particular, the general case
of a sign-changing quartic in three total variables remains unresolved here.

For a number field $K$, write $\mathcal O_K$ for its ring of integers.
If $V\subseteq\mathbb A^n_K$ is an algebraic set with first coordinate $T$,
put
\[
 \pi(V(K))=\{t\in K:\text{some }(t,y_1,\ldots,y_{n-1})\in V(K)\}.
\]
Thus $n$ counts all coordinates, and $n-1$ counts existential witnesses.
All assertions in this supplement are independent of the proposed
arithmetic proof in the repository.

\subsection{Curves and rational parametrizations}

We use the following consequences of Siegel's theorem and the topology
of a number field.  They apply to open subsets, which is important when
using birational maps.

\begin{lemma}\label{p2lb:curves}
Let $V$ be an algebraic set over $K$ of dimension at most one, and let
$f:V\to\mathbb A^1_K$ be a regular function.  If
$f(V(K))\subseteq\mathcal O_K$, then $f(V(K))$ is finite.
The same conclusion holds for a rational function restricted to any
Zariski-open subset on which it is defined.
\end{lemma}
\begin{proof}
There are finitely many irreducible components.  A zero-dimensional
component contributes finitely many points, and a component on which
$f$ is constant contributes at most one value.  On an irreducible curve
which is not geometrically irreducible, every $K$-point lies on the
intersection of two distinct conjugate geometric components.  There
are only finitely many such points.

For a geometrically integral curve, normalization and projective
completion discard or add only finitely many points.  Let $C$ be the
resulting smooth projective curve and let $f\in K(C)$ be nonconstant.
If $C$ has positive genus, Siegel's theorem in its rational-function
form gives only finitely many $P\in C(K)$ for which $f(P)$ is defined
and integral.  One obtains this form from the usual theorem by removing
the poles of $f$ and enlarging a finite set of primes so that the
coordinate functions of that affine curve are integral over
$\mathcal O_{K,S}[f]$; see \cite[Chapter~8]{lbLang1983}.

If $C$ has genus zero and has a $K$-point, it is $K$-isomorphic to
$\mathbb P^1$.  A nonconstant rational function on a nonempty open
subset of $\mathbb P^1(K)$ cannot take only integral values.  Indeed,
embed $K$ diagonally in
$K_\infty=\mathbb R^{r_1}\times\mathbb C^{r_2}$.  The image of $K$
is dense, while that of $\mathcal O_K$ is a discrete lattice.  Continuity
would force the rational function to be constant on a small connected
open ball, and hence identically constant.  If $C(K)$ is empty, only
the previously discarded singular points remain.  These cases prove
the assertion.
\end{proof}

\begin{lemma}\label{p2lb:unirational}
Let $X/K$ be an integral unirational variety, let $U\subseteq X$ be
nonempty and open, and let $f\in K(X)$ be nonconstant and defined on $U$.
Then $f(U(K))$ is not contained in $\mathcal O_K$.
\end{lemma}
\begin{proof}
Compose $f$ with a dominant rational map from projective space to $X$.
After restricting to an affine open set, the composition is a
nonconstant rational function in finitely many variables.  The density
and discreteness argument in the preceding proof, now on
$K_\infty^m$, gives a contradiction if all its values are integral.
\end{proof}

\begin{lemma}\label{p2lb:ruled}
Let $C/K$ be a geometrically integral curve, let $m\geq1$, and let
$X$ be birational over $K$ to $C\times\mathbb P^m$.
Suppose $U\subseteq X$ is nonempty and open and $U(K)$ is Zariski
dense.  A nonconstant $f\in K(X)$, defined on $U$, cannot satisfy
$f(U(K))\subseteq\mathcal O_K$.
\end{lemma}
\begin{proof}
Restrict the birational correspondence to isomorphic nonempty opens
and replace $C$ by its smooth projective model.  We may work on an
open subset of $C\times\mathbb A^m$.  If $f$ is nonconstant in the
affine coordinates over $K(C)$, then outside finitely many points of
$C$ its restriction to a fiber is nonconstant.  Outside another finite
set that fiber meets the chosen open subset in a nonempty open subset
of $\mathbb A^m$.  Density of $U(K)$ supplies a rational base point
outside these exceptional sets.  The restriction to this fiber violates
Lemma~\ref{p2lb:unirational}.

Otherwise $f$ belongs to $K(C)$ and is nonconstant there.  For every
$c\in C(K)$ outside a finite set, the open fiber has a $K$-point,
so $f(c)$ is integral.  Lemma~\ref{p2lb:curves} gives only finitely
many such values, hence only finitely many such $c$, because a
nonconstant map of projective curves has finite fibers.  This
contradicts the density of $U(K)$ in $C\times\mathbb A^m$.
\end{proof}

\begin{remark}
The preceding lemma concerns a variety birational over $K$ to a
product with projective space.  Geometric ruledness after extending
scalars is not sufficient for its proof.  In particular, it must not
be applied without further argument to an arbitrary conic bundle
over $\mathbb P^1_K$ which has no rational section.
\end{remark}

\begin{lemma}\label{p2lb:dense-component}
Let $V/K$ be an algebraic set and $f:V\to\mathbb A^1_K$ a regular
function.  If $f(V(K))$ is infinite and contained in $\mathcal O_K$,
the Zariski closure of $V(K)$ has a geometrically integral component
$X$ of dimension at least two such that $X(K)$ is Zariski dense
and $f|_X$ is nonconstant.
\end{lemma}
\begin{proof}
The closure has finitely many irreducible components, each with dense
$K$-points on the open set obtained by deleting the other components.
If such a component were geometrically reducible, its rational points
would lie in the proper closed intersections of its conjugate
geometric components, a contradiction.  Components on which $f$ is
constant contribute finitely many values, and components of dimension
at most one contribute finitely many by Lemma~\ref{p2lb:curves}.
Thus at least one component has the required properties.
\end{proof}

\subsection{Two quadrics in projective four-space}

The next argument uses only the elementary degree bound
$\deg X\geq1+\operatorname{codim}(X)$ for a nondegenerate projective
integral variety, B\'ezout's theorem, and the following theorem of
Koll\'ar: a geometrically integral cubic hypersurface of dimension at
least two over a perfect field is unirational if it has a rational
point and is not a cone \cite[Theorem~2]{lbKollar2002}.
The elementary degree bound is recalled in
\cite[Introduction]{p2lbEGHP2006}.  We give the projection argument
explicitly in order to include reducible and nonreduced intersections.

\begin{proposition}\label{p2lb:two-quadrics-surface}
Let $Q_1,Q_2\in K[X_0,\ldots,X_4]$ be coprime homogeneous
quadratics, and let $X$ be a geometrically integral surface component
of $V(Q_1,Q_2)$.  Suppose $X(K)$ is Zariski dense.
Then $X$ is either unirational over $K$ or birational over $K$ to
$C\times\mathbb P^1$ for a geometrically integral curve $C/K$.
\end{proposition}
\begin{proof}
The complete intersection $W=V(Q_1,Q_2)\subseteq\mathbb P^4$
has degree four.  More precisely, if $m_X$ is its multiplicity along
$X$, then $m_X\deg X\leq4$.  This follows either from B\'ezout's
theorem with multiplicities or from the complete-intersection Hilbert
series $(1-t^2)^2/(1-t)^5$ and additivity of degree.

If $\deg X\leq2$, the elementary degree bound shows that its linear
span has dimension at most three.  It is therefore a plane or an
integral quadric surface in its span.  Its linear span is defined over
$K$, and density supplies a smooth $K$-point in the quadric case.
Projection from that point makes the quadric rational; a plane is
rational as well.

Suppose $\deg X\geq3$.  The degree inequality forces $m_X=1$.
Since the characteristic is zero, a nonempty open subset of $X$ is
smooth in the complete intersection $W$.  Choose a $K$-point there,
using density, and change coordinates so that it is
$p=(1:0:0:0:0)$.  The quadratics have the form
\[
 Q_1=X_0L_1+q_1,\qquad Q_2=X_0L_2+q_2,
\]
where $L_i$ are linear and $q_i$ are quadratic in $X_1,\ldots,X_4$.
Smoothness of $W$ at $p$ means that $L_1,L_2$ are independent.
The cubic
\[
 H=L_1q_2-L_2q_1
\]
is not identically zero.  Otherwise coprimality of $L_1,L_2$ gives
$q_1=L_1L_3$ and $q_2=L_2L_3$ for a linear form $L_3$, contradicting
coprimality of $Q_1,Q_2$.

The surface $X$ cannot be contained in the plane $L_1=L_2=0$,
because its degree is at least three.  Interchange the two quadratics
if necessary so that $L_1$ does not vanish identically on $X$.
Projection from $p$ then has the rational inverse
\[
 X_0=-q_1/L_1
\]
on the corresponding open set.  Thus $X$ is birational to a
geometrically integral surface component $Y$ of the cubic
$V(H)\subseteq\mathbb P^3$.  The rational points of $Y$ are dense.

If $Y$ is a plane or a quadric, it is rational as above.  If it is
a cubic which is not a cone, Koll\'ar's theorem makes it unirational.
Finally, a geometrically integral cubic surface which is a cone has
a single vertex point defined over $K$.  Indeed the vertex directions
form the kernel of a linear map over $K$ given by the directional
derivatives of its cubic equation; more than one independent vertex
direction would leave a binary cubic, which is geometrically reducible.
The complement of the vertex is birational over $K$ to the product
of its plane cubic base with $\mathbb P^1$.  This proves the result.
\end{proof}

\begin{theorem}\label{p2lb:quadratic-system-four}
Let $q_1,\ldots,q_s\in K[T,Y_1,\ldots,Y_{n-1}]$ have total
degree at most two, where $n\leq4$.  If the rational projection
\[
 S=\{t\in K:\exists\mathbf y\in K^{n-1},\quad
                   q_1(t,\mathbf y)=\cdots=q_s(t,\mathbf y)=0\}
\]
is contained in $\mathcal O_K$, then $S$ is finite.
\end{theorem}
\begin{proof}
Adding unused coordinates reduces the assertion to $n=4$.
Zero residuals can be removed; if all residuals vanish identically,
then $S=K$, contrary to the hypothesis.  A nonzero constant residual
gives the empty set.  Let $V$ be their common zero set.

Suppose $S$ is infinite, and choose $X\subseteq V$ as in
Lemma~\ref{p2lb:dense-component}.  Any common polynomial factor
$g$ of the residuals has degree at most two.  The entire hypersurface
$g=0$ lies in $V$, and its input projection is therefore contained
in $\mathcal O_K$.  The quadratic obstruction proved in the first
report says this projection has at most two elements.  Consequently
$X$ is not contained in $g=0$.  Dividing the residuals by their common
greatest divisor does not remove $X$ and leaves quadratics with
greatest common divisor one.

If their $K$-linear span contains a nonzero affine linear polynomial,
restrict to its affine hyperplane, isomorphic to $\mathbb A^3_K$.
Not all restricted residuals vanish identically, since that would
give the original quadratics a common linear factor.  The resulting
surface $X$ is a hypersurface component of a nonzero quadratic in
three variables, hence is a plane or a quadric.  Density gives a
rational parametrization, contradicting Lemma~\ref{p2lb:unirational}.

Otherwise choose two coprime quadratics in the linear span of the
residuals.  To see that this is possible, first choose a nonzero
quadratic $q$ and then a linear combination not divisible by any
irreducible factor of $q$; finitely many proper linear subspaces
cannot cover a vector space over the infinite field $K$.
Their affine common zero set has dimension two.  Thus $X$ is a
surface component.  Their degree-two homogenizations remain coprime,
and the projective closure of $X$ is a component of their intersection
in $\mathbb P^4$.

Proposition~\ref{p2lb:two-quadrics-surface} makes this component
unirational or birational to a product of a curve and $\mathbb P^1$.
The input is a nonconstant rational function and is integral on
the rational points of its affine open subset.  This contradicts
Lemma~\ref{p2lb:unirational} or Lemma~\ref{p2lb:ruled}.
\end{proof}

\begin{corollary}\label{p2lb:sos-five}
Suppose
\[
 P(T,\mathbf Y)=\sum_{j=1}^s c_jq_j(T,\mathbf Y)^2,
 \qquad c_j\in\mathbb Q_{>0},\quad\deg q_j\leq2.
\]
If the rational existential projection of $P=0$ is $\mathbb Z$,
then $P$ uses at least five total variables, hence at least four
existential witnesses.
The same lower bound holds over a number field for any quartic
obtained by applying an anisotropic quadratic form to quadratic
residuals and defining its ring of integers.
\end{corollary}
\begin{proof}
In the rational case, the sum is zero exactly when all $q_j$ vanish.
For an anisotropic combining form the same equivalence holds by
definition of anisotropy.  Apply Theorem~\ref{p2lb:quadratic-system-four}.
\end{proof}

This improves the unrestricted three-total-variable bound for the
sum-of-squares quartic architecture used by the released polynomial
and all the expanded quartic variants in the first report.  It does
not rule out a sign-changing quartic in three or four total variables.

\subsection{Nonnegative polynomials in three variables}

\begin{proposition}\label{p2lb:nonnegative-squarefree}
Let $P\in\mathbb Q[T,Y,Z]$ be squarefree and nonnegative on
$\mathbb R^3$.  If its rational input projection is contained in
$\mathbb Z$, that projection is finite.
\end{proposition}
\begin{proof}
Every real zero of a differentiable nonnegative function is a local
minimum, so all three partial derivatives of $P$ vanish there.
The singular locus of a squarefree hypersurface in characteristic
zero has codimension at least one in that hypersurface.  In this
case it has dimension at most one.  All rational zeros lie in this
singular locus.  Apply Lemma~\ref{p2lb:curves} to its input function.
\end{proof}

\begin{theorem}\label{p2lb:nonnegative-degree-eight}
Let $P\in\mathbb Q[T,Y,Z]$ be nonnegative on $\mathbb R^3$.
If its rational input projection is an infinite subset of $\mathbb Z$,
then $P$ has an irreducible factor of degree at least four and
multiplicity at least two.  In particular $\deg P\geq8$.
\end{theorem}
\begin{proof}
Choose a component $X$ of the closure of its rational zeros as in
Lemma~\ref{p2lb:dense-component}.  It has dimension two, so its prime
ideal is generated by an irreducible factor $g$ of $P$.  Write
$P=g^m h$ with $g\nmid h$.  Density supplies a rational point on
$X$ which is smooth on $g=0$ and does not lie on $h=0$.
If $m=1$, then $\nabla P=h\nabla g$ is nonzero there, contrary
to nonnegativity of $P$.  Thus $m\geq2$.

The projection of the whole hypersurface $g=0$ is contained in
the projection of $P=0$, and is infinite because the input is
nonconstant on $X$ and $X(\mathbb Q)$ is dense.  The unconditional
degree-three obstruction of the first report gives $\deg g\geq4$.
Therefore $\deg P\geq m\deg g\geq8$.
\end{proof}

Nonnegativity only on rational inputs gives the same result, because
$\mathbb Q^3$ is dense in $\mathbb R^3$ and $P$ is continuous.
The degree-eight conclusion is restricted to nonnegative polynomials;
it is not an unconditional degree-eight lower bound for two witnesses.

\subsection{Further necessary conditions}

\begin{proposition}\label{p2lb:linear-coordinate}
If a polynomial $P(T,Y_1,\ldots,Y_r)\in K[T,\mathbf Y]$ defines
$\mathcal O_K$ in $K$, its degree in each witness which occurs is
at least two.  Its degree in $T$ is also at least two.
\end{proposition}
\begin{proof}
Suppose $P=AY_r+B$, where $A\ne0$ is independent of $Y_r$.
For all but finitely many $t\in K$, the specialization
$A(t,Y_1,\ldots,Y_{r-1})$ is a nonzero polynomial.  Since $K$ is
infinite, choose the other witnesses so that it does not vanish and
then solve uniquely for $Y_r$.  The projection is cofinite in $K$,
which is incompatible with $\mathcal O_K$.

If $P=AT+B$ with $A\ne0$ independent of $T$, its projection includes
the values of $-B/A$ on the nonempty open set $A\ne0$.  Integrality
of all these values forces this rational function to be constant,
by the density and discreteness argument.  Thus $B=-cA$ for some
$c\in K$, and $P=A(T-c)$.  If $A$ has a $K$-zero, every input is
accepted; otherwise only $c$ is accepted.  Neither is $\mathcal O_K$.
If $P$ is independent of $T$, its projection is empty or all of $K$.
\end{proof}

\begin{proposition}\label{p2lb:quartic-surface-necessary}
Suppose a quartic $P\in\mathbb Q[T,Y,Z]$ defines $\mathbb Z$ in
$\mathbb Q$.  Then $P$ is irreducible and geometrically irreducible,
its rational points are Zariski dense on its affine surface, and
$P$ takes both positive and negative values on $\mathbb R^3$.
The input function is nonconstant on that surface.  Every curve
birational over $\mathbb Q$ to $\mathbb P^1$ on which the input is defined is vertical
(the input is constant), and the same holds for every geometrically
integral genus-one curve on the surface with infinitely many rational
points.
\end{proposition}
\begin{proof}
If $P$ factors nontrivially, every factor has degree at most three.
Each factor's rational projection is contained in $\mathbb Z$ and
is finite by the earlier degree-three obstruction.  Their finite
union cannot be $\mathbb Z$.  If the irreducible $P$ is geometrically
reducible, its rational points lie in intersections of distinct
conjugate components, of dimension at most one; Lemma~\ref{p2lb:curves}
again gives a finite projection.  The same lemma shows that a proper
Zariski closure of the rational points would give a finite projection.
Thus the surface is geometrically integral and its rational points
are dense.  The input is nonconstant because its image is infinite.

If $P$ had only one sign, apply
Theorem~\ref{p2lb:nonnegative-degree-eight} to $P$ or $-P$.
Finally, a nonvertical rational curve with a rational point would
violate Lemma~\ref{p2lb:unirational}.  On a genus-one curve with
infinitely many rational points a nonconstant input function would
take only finitely many integral values at rational points by
Siegel's theorem, whereas every rational point where the input is
defined would have to give such a value.  This is impossible.
\end{proof}

\subsection{Why the unrestricted two-witness case remains difficult}

A smooth quartic surface in $\mathbb P^3$ has trivial canonical
bundle and is a K3 surface.  It is not covered by the unirational
cubic theorem.  Zariski density of rational points, even when known,
does not itself imply density in the real topology or the existence
of a nonintegral value of a specified rational function.  Likewise,
potential density after a finite field extension cannot be substituted
for a statement about rational points over $\mathbb Q$.
For examples and the distinction between density over the ground field
and potential density, see \cite{p2lbHarrisTschinkel2000}.

The same limitation appears naturally in the improved bound for
quadratic systems.  A smooth complete intersection of three quadrics
in $\mathbb P^5$ is a K3 surface: adjunction gives
$K_X=\mathcal O_X(2+2+2-6)=\mathcal O_X$, and the
complete-intersection cohomology gives $H^1(X,\mathcal O_X)=0$.
Thus five affine coordinates already permit a class of surfaces not
handled by the two-quadric argument.  This observation is a limit of
the proof, not an existence claim for a five-variable definition.

\begin{proposition}[Conditional obstruction for general-type surfaces]
\label{p2lb:lang-surface}
Assume the weak Bombieri--Lang conjecture for surfaces over $\mathbb Q$:
the rational points of a smooth projective surface of general type are
not Zariski dense.  Then no polynomial in three total variables whose
geometrically integral projective surface has a smooth projective
model of general type can have an infinite rational input projection
contained in $\mathbb Z$.
In particular, under this conjecture a polynomial in three variables
with smooth projective closure of degree at least five cannot define
$\mathbb Z$.
\end{proposition}
\begin{proof}
Lemma~\ref{p2lb:curves} forces Zariski density of rational points
on a surface with such a projection.  Passing to isomorphic dense
opens of a smooth projective model preserves density, contradicting
the conjecture.  A smooth degree-$d$ surface in $\mathbb P^3$ has
canonical bundle $\mathcal O_X(d-4)$, which is ample when $d\geq5$.
\end{proof}

Mazur's conjecture about finitely many real connected components of
the closure of rational points would rule out every Diophantine
definition of $\mathbb Z$ in $\mathbb Q$, in any number of variables.
It therefore cannot supply a conditional numerical improvement
compatible with simultaneously assuming the proposed arithmetic
theorem.  Proposition~\ref{p2lb:lang-surface}, by contrast, is a
restriction on one class of possible two-witness surfaces.

% Includable supplement. Independent of the proposed arithmetic construction.
\section{A restriction on the degree in two witnesses}
\label{sec:p2lb-conics}

The total degree can conceal an important distinction: a polynomial may
have arbitrarily high degree in the input, while remaining quadratic in
its witnesses.  With only two witnesses, such a polynomial cannot define
the ring of integers.  We prove a simultaneous version which also rules
out a finite disjunction of such equations.

\begin{theorem}\label{p2lb:quadratic-witnesses}
Let $K$ be a number field and let
$F_1,\ldots,F_s\in K[T,Y,Z]$ have total degree at most two in the
pair $(Y,Z)$, with no restriction on their degree in $T$.  Put
\[
 S=\{t\in K:\text{for some }i\text{ and }y,z\in K,
                         \ F_i(t,y,z)=0\}.
\]
Exactly one of the following alternatives holds:
\begin{enumerate}
\item at least one equation has a solution
$(Y,Z)\in K(T)^2$, and $S$ is cofinite in $K$;
\item no equation has such a solution, and $K\setminus S$ contains
infinitely many elements of every coset $a+\mathfrak a$, where
$a\in\mathcal O_K$ and $\mathfrak a$ is a nonzero ideal of
$\mathcal O_K$.
\end{enumerate}
In particular, no single polynomial quadratic in two witnesses, or
finite disjunction of such polynomials, defines $\mathcal O_K$ in $K$.
\end{theorem}

The cofinite assertion in the first alternative follows by specializing
the rational functions in a generic solution away from their finitely
many poles.  The substantive issue is to find integral specializations
which simultaneously preserve the obstructions in the second alternative.
We use Faddeev's residue sequence, Chebotarev's theorem and integral
Hilbert irreducibility.  The argument is a consequence of these classical
tools; no claim of historical priority is made.

\subsection{Quaternion obstructions at integral specializations}

For a field $F$ of characteristic different from two, $(a,b)$ denotes the
Brauer class of the quaternion algebra generated by $i,j$ with
$i^2=a$, $j^2=b$ and $ij=-ji$.  A smooth conic is split, or equivalently
has an $F$-point, precisely when its associated quaternion class is zero.
For this correspondence and the residue sequence below see
\cite[Theorem~5.1 and Section~6.1]{p2lbBergEtAl2024}.

We recall explicitly the residue convention.  At the closed point of
$\mathbb A^1_K$ given by a monic irreducible $f(T)$, write
\[
 a=f^r u,\qquad b=f^s v,
\]
where $u,v$ are units in $K[T]_{(f)}$.  If $\theta$ is a root of $f$
and $L=K(\theta)$, the residue of $(a,b)$ is the square class of
\[
             (-1)^{rs}u(\theta)^s/v(\theta)^r
             \quad\text{in }L^\times/L^{\times2}.
\]
Faddeev's theorem implies that a class in $\operatorname{Br}K(T)$
with all finite residues zero comes from $\operatorname{Br}K$.
Indeed the sum-of-corestrictions relation first makes the residue at
infinity zero, since its residue field is $K$; exactness then gives a
constant class.  In characteristic zero the sequence applies to the
whole Brauer group because $\operatorname{Br}\overline K(T)=0$.

We shall also use the following elementary local formula.  If $F$ is a
nonarchimedean local field with odd residue field $\kappa$, if $\varpi$
is a uniformizer, and if $u,v$ are units, the quaternion
$(\varpi^r u,\varpi^s v)$ is split exactly when
\[
                    (-1)^{rs}\overline u^{\,s}/\overline v^{\,r}
                    \quad\text{is a square in }\kappa^	imes.
\]
One obtains this by bilinearity of the Hilbert symbol: two unit
parameters give a split algebra, $(\varpi,u)$ is detected by the square
class of $\overline u$, and $(\varpi,\varpi)=(\varpi,-1)$.
These facts also follow from the norm criterion for the unramified
quadratic extension; see \cite[Proposition~5.5]{p2lbVoight2019}.

\begin{lemma}\label{p2lb:brauer-progression}
Let $\alpha\in\operatorname{Br}K(T)$ be a nonzero quaternion class.
If $\alpha$ is constant, every specialization outside a finite set
is nonzero.  If $\alpha$ is not constant, then outside any prescribed
finite set of primes of $K$ there is a prime $\mathfrak p$ and a
residue class $c\bmod\mathfrak p^2$ such that, outside a finite set
of inputs, every $t\in\mathcal O_K$ in that residue class has
\[
                   \alpha(t)\ne0\quad\text{in }
                   \operatorname{Br}K_{\mathfrak p}.
\]
\end{lemma}
\begin{proof}
The assertion for a constant class follows from specialization on an
open set where its equality with the chosen quaternion presentation is
defined.  Otherwise choose a nonzero finite residue using Faddeev's
theorem.  With the notation just introduced, let
\[
 d=(-1)^{rs}u(\theta)^s/v(\theta)^r\in L^\times
\]
be its nonsquare representative.  Choose a finite Galois extension
$M/K$ containing $L(\sqrt d)$, and an element
$g\in\operatorname{Gal}(M/L)$ which sends $\sqrt d$ to $-\sqrt d$.
Such an element exists because $d$ is not a square in $L$.

Chebotarev's theorem gives infinitely many primes of $K$ with Frobenius
conjugate to $g$; see \cite[Section~2.1, Theorem~1]{p2lbSerre1981}.
Choose one outside the prescribed finite set and all further primes of
bad reduction for the displayed algebraic data, including primes above
two.  Choose a prime of $M$ whose Frobenius is exactly $g$, and let
$\mathfrak q$ be its restriction to $L$.  Since $g$ fixes $L$, this
prime of $L$ has residue degree one over $\mathfrak p$.  Since $g$
acts nontrivially on $\sqrt d$, the reduction of $d$ at
$\mathfrak q$ is a nonsquare.  The excluded bad primes ensure that
it is a nonzero unit.

The corresponding completion of $L$ is $K_{\mathfrak p}$, so
$f$ has a simple root $\theta_{\mathfrak p}$ in
$\mathcal O_{K_{\mathfrak p}}$.  Choose a residue class
$c\bmod\mathfrak p^2$ for which
\[
 v_{\mathfrak p}(t-\theta_{\mathfrak p})=1
 \qquad(t\equiv c\bmod\mathfrak p^2).
\]
There is such a class because the reduction map from $\mathcal O_K$
onto $\mathcal O_{K_{\mathfrak p}}/\mathfrak p^2$ is surjective.
For all these inputs, $v_{\mathfrak p}(f(t))=1$, and $u(t),v(t)$
are units whose reductions equal those at $\theta$ and
$\mathfrak q$.  The local formula therefore makes the specialized
quaternion nonsplit, since its residue square class is that of $d$.
The finitely many inputs where a rational presentation is undefined may
be discarded.  This proves the assertion.
\end{proof}

\begin{corollary}\label{p2lb:simultaneous-conics}
Let $C_1,\ldots,C_m$ be smooth conics over $K(T)$ with no
$K(T)$-point.  Given any coset $a+\mathfrak a$ as in
Theorem~\ref{p2lb:quadratic-witnesses}, it contains a smaller nonempty
coset of a nonzero ideal such that, outside finitely many inputs, none
of the specialized conics has a $K$-point.
\end{corollary}
\begin{proof}
Constant nonzero quaternion classes already exclude every good
specialization.  For each remaining conic, use the preceding lemma
with a fresh prime, distinct from all previously chosen primes and
coprime to $\mathfrak a$.  The Chinese remainder theorem simultaneously
imposes its bad residue class modulo the square of that prime and the
original coset.  Every resulting good specialization is locally insoluble
at the selected prime for each corresponding conic.
\end{proof}

\subsection{Degenerate conics and the simultaneous theorem}

\begin{proof}[Proof of Theorem~\ref{p2lb:quadratic-witnesses}]
Write $F=K(T)$.  If one equation has a point in $F^2$, its rational
functions specialize to a point at every input outside a finite set.
This gives the first alternative.  Suppose henceforth that none does.
We classify each equation over $F$, and discard finitely many inputs
so that all rational coordinate changes, ranks and nonzero coefficients
in this classification specialize correctly.

An equation of witness degree zero is a nonzero element of $F$, and
has no point at any good input.  Witness degree one always gives an
affine $F$-point, so cannot occur in the case under consideration.
For witness degree two, homogenize in $Y,Z$ with a third coordinate
$W$, giving a ternary quadratic form and its projective conic $C$.
If the form has rank three and its quaternion class is zero, then
$C\simeq\mathbb P^1_F$.  The affine open $W\ne0$ is nonempty
and hence has an $F$-point, a contradiction.  Thus every rank-three
case has a nonzero quaternion class, and
Corollary~\ref{p2lb:simultaneous-conics} simultaneously excludes all
of them on a smaller coset of the given integral progression.

A rank-one ternary form is a scalar times the square of an
$F$-linear form.  That line cannot be the line at infinity, since the
original polynomial has witness degree two.  It therefore has an
affine $F$-point, again a contradiction.  A split rank-two form is the
product of two distinct $F$-linear forms.  At least one line is not
the line at infinity, and supplies the same contradiction.

The remaining possibility is a nonsplit rank-two form.  After a linear
change over $F$ it is a nonzero scalar times
\[
                         U^2-d(T)V^2,
                         \qquad d(T)\notin F^{\times2}.
\]
Its only projective $F$-point is its radical point, defined by
$U=V=0$.  If that point were affine, it would give an $F$-point of
the original equation.  Therefore its original $W$ coordinate is
zero.  At every good specialization for which $d(t)$ remains a
nonsquare, the same reasoning leaves only this radical point, still
at infinity, and the specialized affine equation has no solution.

There are only finitely many nonsquare functions $d(T)$ to treat.
The smaller coset already obtained contains a progression
$t=c+bz$ with $c\in\mathcal O_K$, $0\ne b\in\mathcal O_K$
and $z\in\mathcal O_K$: choose $b$ in its defining nonzero ideal.
The polynomials
\[
                         Y^2-d(c+bZ)
\]
remain irreducible over $K(Z)$, because $T\mapsto c+bZ$ is an
isomorphism of rational function fields.  The integral Hilbert
irreducibility theorem, simultaneously for this finite list, gives
infinitely many integral $z$ for which all these specializations
remain irreducible and all previously excluded finite conditions are
avoided.  See the integral form recalled in the original lower-bound
section and \cite[Introduction]{lbMuller2001}.  If the list is empty,
any infinitely many good inputs in the coset suffice.

The resulting inputs lie in the original coset $a+\mathfrak a$,
and none admits a solution of any equation.  This proves the second
alternative and the theorem.
\end{proof}

\begin{corollary}\label{p2lb:two-witness-degree-three}
If $P(T,Y,Z)$ defines $\mathbb Z$ in $\mathbb Q$, then its total
degree in the two witnesses $Y,Z$ is at least three.  More generally,
after factoring over $\mathbb Q(T)$, at least one factor has total
witness degree at least three.  This remains true after every
invertible affine change of the witness coordinates over
$\mathbb Q(T)$, provided the transformed equation and inverse change
are defined at all but finitely many input values.
\end{corollary}
\begin{proof}
The first assertion is the theorem with one equation.  If all factors
had witness degree at most two, clear their rational-function
denominators.  Outside finitely many inputs, the original zero condition
is their finite disjunction, contradicting the theorem.  An invertible
affine change over $\mathbb Q(T)$ preserves the rational solution sets
at every good specialization, so the same argument applies after it.
\end{proof}

\begin{corollary}\label{p2lb:conic-integral-superset}
Let $S\subset K$ contain $\mathcal O_K$, and suppose that $K\setminus S$
is infinite.  Then $S$ is not a finite union of projections of equations
in two witnesses which are jointly quadratic in those witnesses, even if
the degree in the input is unrestricted.
\end{corollary}
\begin{proof}
The second alternative of Theorem~\ref{p2lb:quadratic-witnesses} would
omit infinitely many elements of $\mathcal O_K$, while the first would
make $K\setminus S$ finite.  Both contradict the stated assumptions.
\end{proof}

For example, take
\[
 S=\mathbb Z_{(2)}\cap\mathbb Z_{(3)}
   =\{a/b\in\mathbb Q:\gcd(a,b)=1,\ b>0,\ 2\nmid b,\ 3\nmid b\}.
\]
This set contains $\mathbb Z$ and has infinite complement, for instance
the distinct elements $2^{-n}$ for $n\ge1$.  Thus a quadratic
three-witness definition of this semilocal ring cannot be replaced by a
quadratic two-witness definition, or a finite union of such definitions.
The obstruction concerns joint witness degree; it allows arbitrary input
degree and does not assert a lower bound for equations of higher witness
degree.

This is an unconditional restriction on a possible two-witness
definition at every input degree.  It does not exclude general quartic
surfaces, whose generic fibers over the input can have positive genus.
It also does not extend by this proof to quadratic equations in three
witnesses: their fibers are quadric surfaces, and are not controlled by
a single quaternion obstruction of the kind used above.

\bibliographystyle{alpha}
\bibliography{../../optimization/paper_parts/lower_bounds,references}
\end{document}
